% 第46章分段实写草稿；未完成前不得替换正式章节或计入完成。
\chapter{运行相变：探索、执行与结构更新的临界控制}
\label{chp:phase_transition_meta}

智能系统不会始终以同一种方式运行：熟练任务需要稳定而低分支的执行，未知问题需要扩展候选和主动试探，持续失配又可能要求修改结构。把这些变化统称为温度、熵或临界性虽然直观，却容易把采样随机性、表示结构、控制增益、资源负载和物理能耗混成一个量。本章把“相变”限定为可观测运行统计随控制参数变化而发生的定性重组。我们先定义任务相关序参量，再区分稳定执行、无序搜索和自适应临界区间；随后讨论宏观控制器如何调节探索并提交结构更新，为什么LLM末端采样温度不是系统温度，以及何种证据才足以支持涌现判断。物理相变提供模型工具而不是机制证明；HSF-HD给出一般状态动力学，AgentOS中的SPC、TCE与CCM只是其中一种可执行实现。

\section{任务相关序参量：结构复杂度、传播范围与目标对齐}
\label{sec:section_6298}

序参量的作用，是把无法逐项追踪的高维运行状态压缩成可比较的宏观读出，而不是量化所谓“灵魂”。设时刻 $t$ 的联合状态为 $z_t=(h_t,s_t,b_t,r_t)$：$h_t\in\mathbb R^d$ 是快速活动状态，$s_t=(G_t,\Theta_t)$ 是带版本的局部结构网络，$b_t$ 是对环境和候选策略的全局分布表示，$r_t$ 是资源账本。任务规范 $q_t$ 给出查询族、时间窗、容许误差与硬约束。我们定义序参量向量
\[
 m_t=P_{q_t,W}(z_{t-W:t})=(S_t,C_t,A_t,L_t,R_t),
\]
其中 $W$ 是观测窗；$S_t$ 描述任务相关结构复杂度，$C_t$ 描述扰动传播与恢复，$A_t$ 描述动作对目标的对齐，$L_t$ 描述负载和拥塞，$R_t$ 描述结构周转。投影 $P$ 的任务、窗口、归一尺度和数据来源都是定义的一部分。同一系统换任务或换窗口时，序参量可改变，因此不存在脱离测量协议的普遍“智能温度计”。

旧稿的图谱熵密度可以保留为结构读出，但需要修正其解释。对固定时间窗内的无向加权图 $G_t=(V_t,E_t,W_t)$，令归一化拉普拉斯为 $L_t=I-D_t^{-1/2}W_tD_t^{-1/2}$，非负特征值为 $\lambda_i$。若 $\sum_i\lambda_i>0$，定义 $p_i=\lambda_i/\sum_j\lambda_j$ 与
\[
 S_{spec}(G_t)=-\frac{1}{\log n_t}\sum_{i=1}^{n_t}p_i\log p_i.
\]
它无量纲且落在声明范围内，但只描述所选图、权重和节点集的谱分散。低值不必然意味着偏执或晶体，高值也不必然意味着随机、幻觉或精神病；不同非同构图甚至可能同谱。若节点数和权重随时间变化，必须先定义对齐、重采样或谱测度距离。真正与任务有关的结构复杂度还应加入查询保持项，例如在边遮蔽、节点移除和权重扰动下，读出 $Q_q$ 的变化 $\Delta_q=\sup_{\delta\in\mathcal D}d(Q_q(G),Q_q(G+\delta))$。这样才能区分“谱很丰富但任务脆弱”和“谱较集中但任务稳健”。

旧稿的场连贯性 $\xi$ 可重建为传播范围和恢复尺度。令 $a_i(t)$ 是经身份对齐的节点活动，先扣除共同输入和基线，计算条件相关函数 $C_q(i,j;\tau)$。在固定图距离 $d_G$ 和局部平稳区间内，若经验包络可由 $c_0d^{-\alpha}\exp(-d/\xi)$ 合理拟合，则 $\xi$ 是该模型下的相关长度；拟合失败时不应强报一个数。相关不等于因果：共同广播可让远端节点同步，却不表示信息沿路径传播。我们因此补充扰动响应
\[
 C_t^{int}(r)=\mathbb E\left[d\bigl(Q_q(z_{t+\tau}^{do(i,\epsilon)}),Q_q(z_{t+\tau})\bigr)\mid d_G(i,j)=r\right]
\]
以及恢复时间 $\tau_{rec}$。传播范围变长既可能提高跨域整合，也可能放大错误级联；无限相关长度不是顿悟的定义，更不能由一次全局同步推出宏观量子态。

旧稿的规范极化度 $M$ 可改写为目标对齐度。对窗口内实际提交的动作 $u_k$，定义目标方向 $g_k=-\nabla_uJ_{q_k}$ 或离散可行改进集合，令
\[
 A_t=\frac{\sum_{k=t-W+1}^{t}w_k\langle u_k,g_k\rangle}
 {\sum_{k=t-W+1}^{t}w_k(\|u_k\|\|g_k\|+\epsilon)}.
\]
$A_t$只在动作与目标方向位于同一内积空间时成立。高对齐表示动作一致地追随当前目标，不表示目标正确；极端高值可能是刚性服从，低值可能来自多目标权衡、探索动作、执行器故障或目标版本冲突。故还要记录约束违例率、策略多样性、外部任务成绩和目标来源。企业所有员工都对齐KPI可能同时集体刷指标，这正是“意志统一”等于智能的反例。

结构复杂度、传播范围和目标对齐不能合成一个未经校准的总分。我们保留向量序参量，并在任务需要时用Pareto区域或带硬约束的分层判定。局部结构网络负责节点身份、关系、来源和权限；全局分布表示负责候选、置信度与不确定性；二者通过身份绑定、消息耦合和读出接口联合。AgentOS可以由SPC压缩 $h(t)$ 的任务状态，由CCM估计负载、循环和传播范围，由TCE调节控制增益；但一般智能系统可以采用其他观测器。序参量必须能由日志重算、随版本复现并对干预作出可预期响应，否则它只是富有感染力的命名。

从统计上识别运行区间，需要先建立基线而非凭单点阈值。我们在多个任务、随机种子和扰动强度上采集 $m_t$，估计条件分布 $p(m\mid q,c)$，再用变化点、迟滞环、双峰分布或临界减速等证据判断是否存在状态重组。仅看到幂律直线或熵值居中，不能证明自组织临界；有限样本、混合分布和测量截断都能制造相似外观。序参量在本章因此是可证伪的运行摘要：它帮助定位何时换控制策略，却不替代因果模型、外部验收和安全边界。

\section{运行相态：稳定执行、无序搜索与自适应临界区间}
\label{sec:section_6651}

我们把“晶体、气体与液体”保留为三类运行区间的辅助名称，不把它们当作智能载体的真实物态。系统的控制参数记为 $c_t=(\tau_t,g_t,\rho_t,b_t^{res})$：$\tau_t$ 是候选采样尺度，$g_t$ 是目标与抑制增益，$\rho_t$ 是结构更新率，$b_t^{res}$ 是资源预算。快状态 obey $h_{k+1}=F(h_k,s_k,o_k,u_k,\eta_k;c_k)$，慢结构由事务映射 $s_{k+1}=U(s_k,\mathcal E_k;c_k)$ 更新。相态是统计量 $m_t$ 在给定任务族中的稳定分区；当参数连续变化而占用分布、恢复时间或结构周转发生显著跃迁和迟滞时，我们才使用相变一词。

稳定执行区间对应旧稿的层流或晶体态。其典型特征是候选有效数低、策略切换少、目标对齐高、扰动增益有界、结构版本基本不变。对已经校准的算术、编译、熟练驾驶或重复工具调用，这种区间能降低时延并提高复现性。若固定结构和目标下存在Lyapunov函数 $V$，使
\[
 V(h_{k+1})-V(h_k)\le-\alpha\|e_k\|^2+\gamma\|w_k\|^2,
\]
则可得到局部输入到状态界，但这不等于“耗散最低”或零能耗。执行仍消耗计算资源和电能，稳定也不保证输出事实正确。其失效模式是刚性：面对分布外观测、目标变化或不可行约束时，系统继续复用旧轨迹。老司机在熟悉道路上高效，遇到道路塌方若仍沿旧路线行驶，就是稳定控制超出适用域。

无序搜索区间对应旧稿的气体或离子态。这里候选有效数、状态总变差和策略切换率升高，目标对齐下降，验证积压与资源占用上升。它不必是病理：头脑风暴、广域检索和故障排查早期都需要暂时扩大搜索。危险来自搜索与验收脱耦——生成速率超过验证速率时，候选在 $h(t)$ 中拥塞，来源绑定丢失，局部错误经共享通道级联。我们用无量纲负载
\[
 L_t=\frac{\lambda_{in}(t)\,\mathbb E[c_{verify}(t)]}{\mu_{verify}(t)+\epsilon}
\]
表示单位窗内输入到达和验证服务能力的比值。$L_t>1$持续存在提示队列发散风险，但它不是雷诺数，也不推出人的临床状态。对LLM，高采样温度可能增加词元多样性，却不必增加内部推理搜索，更不能把幻觉单独归因于“高温”。

自适应临界区间位于可重复执行与广域搜索之间：系统保持足够候选以逃离局部策略，又能用身份、约束和回执迅速筛除无效分支。其操作判据不是“关联长度无穷大”，而是在预算内同时满足四项条件：对小扰动有可测但有界的响应；任务收益对适度探索呈正增益；错误级联可在规定时间恢复；结构更新经保留测试后提高跨情境表现。我们可定义任务敏感度 $\chi_q=\partial\mathbb E[Y_q]/\partial a$，其中 $a$ 是受控扰动幅度；只有在明确区间内估计，不能外推到奇异发散。真正的系统若响应无界，往往意味着不稳定而非创造力。

洗澡时突然想到数学解的旧例，可以解释为后台重组、线索变化和候选读出的共同结果，不需要假设全脑瞬时锁相或调和零模。要检验“临界区间有利于洞见”，应预注册任务、控制练习时间和休息，记录候选数量、验证链、解决率与恢复时间，并与固定低探索和高探索基线比较。一次戏剧性顿悟不能证明相变；反过来，渐进式证明也同样是创造性成果。物理临界现象提供有关响应、尺度和涨落的数学语言，但人类经验只构成启发性证据。

相态之间通常存在迟滞和路径依赖。升高探索尺度从稳定执行进入搜索的阈值，不一定等于降低探索重新回到执行的阈值；结构在高探索期间被修改后，返回路径更会改变。我们以混合系统表示：连续状态在模式 $m\in\{execute,search,adapt\}$ 内演化，守卫条件触发模式切换，重置映射处理缓存、积分器、候选集和权限。每次切换必须记录原因、前后参数、未决动作和回滚点。否则频繁“调温”会产生抖振，系统在探索与执行间来回切换，反而浪费预算。

图像 `figure/p_t.pdf` 仍应保留，但图中的轴必须解释为运行控制量而非真实压强和热力学温度。横轴可表示探索尺度或候选熵，纵轴可表示目标控制增益，区域边界由实验估计并带置信区间；不同任务和载体会得到不同边界。SPC可提供任务状态压缩，TCE可在相图上选择模式，CCM监视负载与恢复，Boundary限制危险探索，外部记忆 $M_e$ 保存跨期证据。这套AgentOS实现展示如何落地，却不把三相运行区间规定为所有智能系统的唯一分类。

% 精确续写位置：下一段从旧标签 sec:section_54079 开始，重写“相变控制机制”，随后依次完成 LLM 温度诊断与涌现判据。正式章节保持 queued。
